\documentclass[12pt,a4paper]{article}
\usepackage[margin=2.5cm]{geometry}
\usepackage{parskip}
\usepackage{hyperref}
\usepackage{amsmath}
\usepackage{booktabs}
\usepackage{graphicx}
\usepackage{setspace}
\onehalfspacing

\title{\textbf{Antimicrobial Resistance and Climate: A Network Science Approach}\\
\large Personal Report --- Network Science (NS in Paris)}
\author{Clark Rodriguez}
\date{May 2026}

\begin{document}
\maketitle

% ============================================================
\section{What I Did}
% ============================================================

This project asked whether countries with hotter climates share more
similar antimicrobial resistance (AMR) profiles---and, if they do,
whether that similarity organises them differently within a
country-to-country similarity network.

I built the entire analysis pipeline from scratch in Python, working
through eight sequential steps.
\begin{enumerate}
  \item I downloaded HadEX3 climate indices (TXx: the annual hottest-day
    temperature) in NetCDF format for 1950--2018 and wrote a script to
    extract country-level annual means for 27 EU/EEA countries over 2000--2018,
    using \texttt{regionmask} for country masking and cosine-latitude weighting
    for area-correct spatial averaging.
  \item I downloaded and cleaned ten ECDC AMR surveillance files covering
    percentage resistance rates for five pathogen--antibiotic combinations
    across EU/EEA countries. I filtered to hospital-acquired isolates only,
    harmonised country codes (e.g.\ ECDC uses \texttt{EL} for Greece and
    \texttt{UK} for the United Kingdom), and merged everything into a single
    master dataset of 3\,926 rows across 27 countries.
  \item I split countries at the median TXx (33.66\,\textdegree{}C) into a
    LOW-temperature group (13 countries) and a HIGH-temperature group
    (14 countries). I then built a $27 \times 10$ resistance matrix and
    computed pairwise cosine similarity between every pair of countries.
    I iterated on the similarity threshold---testing 0.90, 0.97, 0.93, 0.95,
    and finally settling on \textbf{0.96}---because that threshold produced a
    fully connected network of 66 edges with zero isolated nodes and a density
    of 0.188.
  \item I computed four node-level centrality metrics (degree, betweenness,
    closeness, and eigenvector) and ran Mann--Whitney U tests to compare the
    LOW and HIGH groups on each metric and on mean resistance rate.
  \item I generated a publication-quality network figure and a set of
    diagnostic plots (centrality boxplots, temperature-vs-resistance scatter
    with OLS regression).
  \item I ran Louvain community detection on the final network, producing four
    communities with a modularity of 0.506.
  \item I wrote a \texttt{.gitignore}, removed raw data from version control,
    and documented the full pipeline in handoff materials for my two junior
    collaborators.
\end{enumerate}

% ============================================================
\section{What the Network Represents}
% ============================================================

Each node in the network is one EU/EEA country.
Two countries are connected by an edge if their resistance profiles
(the vector of mean percentage resistance across 10 pathogen--antibiotic
combinations) are at least 96\% similar by cosine similarity.
The weight of each edge is the exact cosine similarity value.

A high cosine similarity between two countries means they have
proportionally equivalent resistance patterns: if Country A has high
resistance to \textit{E.\ coli}--fluoroquinolones and low resistance to
\textit{K.\ pneumoniae}--carbapenems, Country B has the same shape.
It does not mean that both countries have identical absolute resistance
rates---only the pattern (direction in 10-dimensional feature space) is
compared.

The threshold of 0.96 is high by design.
I wanted edges to represent genuinely strong profile similarity, not just
moderate correlation.
At 0.96, 66 of the 351 possible country pairs are connected; at 0.90,
the graph was so dense (almost complete) that no meaningful community
structure emerged.

Node size in the figures encodes degree centrality (for the publication
figure) or betweenness centrality (for the analysis map), both of which
measure how central a country is to the similarity network.
Node colour encodes temperature group: blue for LOW, red for HIGH.

% ============================================================
\section{Key Findings I Can Explain}
% ============================================================

\textbf{Finding 1: HIGH-temperature countries have significantly higher
mean AMR resistance rates.}

The Mann--Whitney U test on mean resistance rate returned
$U = 26.0$, $p = 0.0017$.
LOW-temperature countries had a median resistance rate of 15.4\%,
while HIGH-temperature countries had a median of 29.5\%---a difference
of more than 14 percentage points.
The linear regression of country mean TXx against mean resistance rate
gave $R^2 = 0.504$, with a slope of $1.68\%$ per \textdegree{}C.
In plain language: for each extra degree of hottest-day temperature,
a country's average AMR resistance rate is associated with an increase
of roughly 1.7 percentage points.
This relationship is statistically significant and explains about half
of the variance across countries.

\textbf{Finding 2: LOW-temperature countries have higher eigenvector
centrality.}

The Mann--Whitney U test on eigenvector centrality returned
$U = 133.0$, $p = 0.044$.
LOW-temperature countries had a median eigenvector centrality of 0.232
versus 0.0002 for HIGH-temperature countries.

Eigenvector centrality measures how well-connected a node's neighbours
are.
A high value means a country is connected to other well-connected
countries---i.e., it sits in a dense, mutually reinforcing hub.
Germany (DEU, $0.412$), Austria (AUT, $0.398$), the Netherlands (NLD,
$0.346$), and Denmark (DNK, $0.316$) are all LOW-temperature countries
with high eigenvector centrality.
They form a tightly interconnected northwestern European cluster whose
resistance profiles are mutually similar.

The interpretation is somewhat counterintuitive: HIGH-temperature
countries have higher absolute resistance rates, but they are not
forming a tight eigenvector-connected hub---they are scattered across
different parts of the network.
LOW-temperature countries, despite lower absolute resistance, are more
structurally cohesive in the similarity space.

\textbf{Finding 3: Four communities, partially stratified by temperature.}

The Louvain algorithm identified four communities (modularity $= 0.506$):

\begin{itemize}
  \item \textbf{Community 0} (11 countries): 8 LOW + 3 HIGH---dominated
    by northern and western Europe (Norway, Ireland, Sweden, Finland,
    UK, Denmark, Netherlands, Germany, plus Belgium, Austria, France).
  \item \textbf{Community 1} (2 countries): both HIGH---Czechia and
    Slovakia.
  \item \textbf{Community 2} (8 countries): 4 LOW + 4 HIGH---a mixed
    central and eastern European cluster (Latvia, Lithuania, Poland,
    Croatia, Italy, Romania, Bulgaria, Greece).
  \item \textbf{Community 3} (6 countries): 2 LOW + 4 HIGH---Iceland,
    Estonia, Slovenia, Hungary, Spain, Portugal.
\end{itemize}

Communities 0 is clearly geographically and climatically coherent
(low-temperature, western Europe).
Community 1 is the smallest and most homogeneous---two neighbouring
countries with nearly identical resistance profiles.
Communities 2 and 3 are more mixed, suggesting that shared resistance
patterns in those clusters are driven by factors beyond temperature
alone (possibly shared antibiotic stewardship policies, healthcare
system links, or travel patterns).

% ============================================================
\section{Limitations I Am Aware Of}
% ============================================================

\textbf{Ecological fallacy.}
All analysis is at country level.
A country's mean resistance rate and mean TXx are aggregated from
individual clinical isolates and grid-cell temperatures respectively.
The correlation I find does not imply that individual patients in hotter
regions are more likely to carry resistant bacteria---it only describes
a country-level pattern.

\textbf{Confounders.}
The regression $R^2 = 0.504$ means that temperature explains about half
the variance in resistance rates across countries.
The other half is unaccounted for.
Likely confounders include antibiotic consumption per capita, healthcare
system capacity, population density, sanitation infrastructure, and
international travel volume.
I collected some of these (GDP per capita, population density,
sanitation, water access) in the master dataset but did not use them
in the network or regression; a multivariate model would be needed to
isolate the temperature signal.

\textbf{Cosine similarity is not epidemiological proximity.}
The network edges represent profile shape similarity, not
transmission or travel links.
Two countries can be close in the network simply because they happen
to have selected for resistance to the same drug classes, not because
they are exchanging resistant strains.

\textbf{Betweenness centrality and edge weights.}
NetworkX's \texttt{betweenness\_centrality} with \texttt{weight} treats
weights as distances (lower = shorter path), but my edge weights are
cosine similarities (higher = closer).
The resulting betweenness values are counter-intuitive: Poland (POL)
tops the list at 0.532 despite being a mid-range country by resistance
profile.
A correct implementation would invert the weights
($d = 1 - \text{similarity}$) before computing betweenness.
This bug does not affect eigenvector or degree centrality, which
were the two metrics that entered the final statistical comparisons.

\textbf{Threshold choice.}
The threshold 0.96 was chosen to produce a connected graph with fewer
than 80 edges, which are reasonable structural constraints for a
network of 27 nodes.
However, the threshold is arbitrary; results would differ at 0.95
(96 edges) or 0.97 (fragmented).
The qualitative conclusion---that HIGH-temperature countries have
higher resistance rates---is robust to the threshold choice because it
comes from the Mann--Whitney test on raw resistance rates, not from
network structure.

\textbf{Data coverage.}
ECDC surveillance covers 30 species--antibiotic combinations but only
10 had sufficient country coverage for the network.
Countries and years with $>$50\% missing combinations were dropped before
imputation, which reduced the dataset from 30 potential countries to 27.

% ============================================================
\section{My Role in the Group}
% ============================================================

I was the network science lead.
I designed the analysis pipeline, chose the methodological approach
(cosine similarity network, Louvain community detection, Mann--Whitney
testing), wrote all eight Python scripts, debugged the data extraction
and preprocessing issues, and generated all figures and outputs.

Marc Jacob Doria and Gerardo Luis Fernando are genomics juniors on the
team.
Their primary background is in microbiology and genomic epidemiology,
not network analysis.
My role included preparing a separate technical handoff document
(\texttt{junior\_handoff.tex}) that explains the pipeline outputs
in biological terms, provides ready-to-paste Results paragraphs for
any written deliverable they produce, and documents the full
country-level data tables so they can answer reviewer questions
without needing to re-run the code.

The split of responsibilities reflects our complementary expertise:
I handled the quantitative and computational side; they contributed
domain knowledge about resistance mechanisms, clinical interpretation,
and the microbiological significance of the pathogen--antibiotic
combinations selected by ECDC surveillance.

\end{document}
